\section{胶子螺旋度分布}\label{sec:helicity}

按照同样的步骤，我们可以定义三个可用于计算胶子螺旋度分布的算符：
\begin{align}
\Delta {O}^{1}_R(z_2, z_1) &= Z_{11}^2 e^{\overline{\delta { m}}\Delta z}  \epsilon_{ij}  F^{ti}(z_2) L(z_2,z_1) F^{tj}(z_1), \non\\
 \Delta{  O}^{2}_R(z_2,z_1) &=  Z_{22}^2 e^{\overline{\delta { m}}\Delta z} \epsilon_{ij}  F^{zi}(z_2) L(z_2,z_1) F^{zj}(z_1),\non\\
\Delta {O}^{3}_R(z_2,z_1) &=  Z_{11}Z_{22} e^{\overline{\delta { m}}\Delta z}  \epsilon_{ij} F^{ti}(z_2) L(z_2,z_1) F^{zj}(z_1),
\end{align}
其中 $\epsilon_{ij}$ 是二维反对称张量。
